\documentclass[11pt]{article}

% ------------------------------------------------
% Page Geometry
% ------------------------------------------------
\usepackage[margin=1in]{geometry}

% ------------------------------------------------
% Fonts and Microtypography
% ------------------------------------------------
\usepackage[T1]{fontenc}
\usepackage{lmodern}
\usepackage{microtype}

% ------------------------------------------------
% Mathematics
% ------------------------------------------------
\usepackage{amsmath, amssymb, amsthm, mathtools}
\usepackage{bm}

% ------------------------------------------------
% Theorem Environments
% ------------------------------------------------
\newtheorem{theorem}{Theorem}
\newtheorem{corollary}{Corollary}
\newtheorem{proposition}{Proposition}
\newtheorem{definition}{Definition}
\newtheorem{remark}{Remark}

% ------------------------------------------------
% Tables and Figures
% ------------------------------------------------
\usepackage{graphicx}
\usepackage{booktabs}
\usepackage{array}
\usepackage{tabularx}
\usepackage{caption}

% ------------------------------------------------
% Hyperlinks
% ------------------------------------------------
\usepackage[hidelinks]{hyperref}

% ------------------------------------------------
% Custom Commands
% ------------------------------------------------
\newcommand{\RID}{\text{RID}}
\newcommand{\Tman}{\mathcal{T}}
\newcommand{\Ggen}{\mathcal{G}}
\newcommand{\Kkey}{\mathcal{K}}
\newcommand{\Oobs}{\mathcal{O}}

% ------------------------------------------------
% Title Information
% ------------------------------------------------
\title{Hyperpleonastic Legibility:\\
Selector Fields, Recoverable Information Density,\\
and Civilizational Incompressibility}

\author{Flyxion}

\date{\today}

\begin{document}

\maketitle

\begin{abstract}

Contemporary alignment discourse focuses primarily on aligning the objectives of increasingly capable optimization systems with human preferences. This paper proposes a complementary geometric constraint: civilizational stability depends on preserving recoverable interpretive structure within the material substrate of everyday artifacts. 

We introduce the concepts of the tolerance manifold, selector fields, and Recoverable Information Density (RID) to formalize how bounded manufacturing variation can encode template-selection logic without degrading functionality. We show that increasing RID enlarges the effective dimensional control surface available to human agents, thereby improving the stability margin of a Nyquist-style containment inequality under temporal and capability asymmetry.

We define Dimensional Poverty as the contraction of civic control surfaces due to opacity and externalized generative grammars. Through case studies ranging from thermovariable textiles to domestic appliances, we demonstrate how hyperpleonastic design — embedding interpretive authority into geometry itself — produces distributed damping and graceful degradation under infrastructural stress.

Alignment, in this framework, is not solely a software safety problem. It is a material constraint on compression: optimization must not collapse tolerance manifolds below the threshold required for distributed legibility. Civilizational incompressibility is therefore a stability condition for a world of accelerating optimization.
\end{abstract}

\newpage
\section{Introduction}

The problem of alignment is typically framed as a mismatch between the objectives of advanced optimization systems and the values of the humans they affect. As computational capability increases and automated systems operate at ever shorter action cycles, oversight becomes temporally and dimensionally strained. Much work in artificial intelligence safety has therefore focused on value specification, objective robustness, and adversarial containment.

This paper addresses a related but distinct structural variable: the dimensional control surface available to human agents interacting with optimized systems. Even if objectives are nominally aligned, oversight can fail when interpretive authority collapses into centralized, opaque subsystems. In such regimes, artifacts become functionally dependent on external databases, firmware locks, and revocable keys. The material world becomes mute.

We propose that alignment must also be understood as a geometric constraint on compression. Every engineered artifact possesses a tolerance manifold — a bounded region of permissible variation that does not degrade function. Conventional optimization seeks to minimize this variation for efficiency and uniformity. We argue that such slack constitutes latent channel capacity. When mapped through a selector field into a public template grammar, it becomes a distributed interpretive substrate.

To quantify this phenomenon, we introduce Recoverable Information Density (RID), defined as the fraction of an artifact’s generative grammar inferable from its physical substrate under bounded observation. RID expands the effective dimensional control surface available to users and thereby increases the stability margin of oversight under capability and temporal asymmetry.

We then define Dimensional Poverty as the civic condition in which recoverable information across critical artifacts falls below a threshold necessary for local repair and substitution. Dimensional Poverty is not economic scarcity but interpretive contraction. It manifests when households and institutions depend on opaque systems whose generative grammars are externalized and revocable.

Through case studies in wearable thermal membranes, domestic fermentation devices, and everyday materials, we demonstrate how hyperpleonastic design — embedding template-selection logic into geometry itself — preserves legibility under degradation and reduces centralized chokepoints. Under crisis conditions, selector-field artifacts degrade gracefully rather than catastrophically.

The thesis advanced here is modest but structural: civilizational stability in an era of accelerating optimization depends not only on aligning objectives, but on preserving distributed legibility. Optimization must not collapse tolerance manifolds below the threshold required for recoverable interpretation. Incompressibility, understood as the persistence of interpretive structure within matter, is therefore a precondition for durable alignment.

\section{Compressive Capture and the Externalization of Grammar}

Before introducing selector fields, we must examine the dominant
structural tendency of contemporary industrial systems:
the externalization of generative grammar.

\subsection{Artifacts and External Keys}

Let an artifact $A$ be produced by a generative grammar $\mathcal{G}$.
In principle, $\mathcal{G}$ determines the artifact’s material composition,
assembly logic, operational constraints, repair pathways, and substitution
limits. The grammar specifies not only what the artifact is made of,
but how it is constructed, how it functions within tolerances, how it
fails, and how it may be restored or adapted.

In traditional craft regimes, substantial portions of $\mathcal{G}$ were
inferable from inspection. Joinery revealed assembly sequences; fiber
orientation indicated load-bearing direction; fastener geometry exposed
disassembly order. The artifact carried traces of its own generative
history within its visible structure. Interpretation required skill,
but it did not require proprietary authorization.

In contemporary industrial systems, by contrast, this grammar is
increasingly externalized. The object persists physically, but the
logic governing its composition and repair is displaced into manuals,
firmware layers, diagnostic software, or remote databases. 

Let $\mathcal{K}$ denote the interpretive key required to access
$\mathcal{G}$. In opaque systems we have

\[
\mathcal{K} \not\subset A,
\]

meaning that the information necessary to reconstruct the artifact’s
generative grammar is not materially embedded within the artifact itself.
Instead, the key resides in proprietary documentation, firmware
authentication layers, remote servers, or licensed diagnostic tools.
Access to interpretation therefore depends on institutional permission
rather than physical inspection.

The artifact remains physically present, but its interpretive authority
is displaced. It functions, yet it does not explain itself; the capacity
to decode its grammar has been separated from its material form.

\subsection{Compression as Design Imperative}

Optimization pressure favors uniformity.
Variance is minimized.
Surface irregularities are polished away.
Manufacturing tolerances are tightened toward canonical form.

Formally, the tolerance manifold $\mathcal{T}$ is contracted:

\[
\dim(\mathcal{T}) \downarrow.
\]

As $\mathcal{T}$ contracts, latent channel capacity collapses and
Recoverable Information Density correspondingly declines. The tolerance
manifold no longer carries surplus variation capable of encoding
interpretive structure; it approaches canonical form.

This contraction is not incidental but structurally incentivized.
Uniformity simplifies mass production, secrecy protects intellectual
property, and centralization concentrates value within institutional
boundaries. Each of these pressures rewards the minimization of
variation and the externalization of grammar.

The side effect, however, is a reduction in dimensional control at the
user level. As slack disappears, so too does the capacity for local
interpretation, adaptation, and reconstruction.

\subsection{Compressive Capture}

We define \emph{Compressive Capture} as the condition in which:

\[
I_{\text{recoverable}}(A) \rightarrow 0
\quad \text{while} \quad
I_{\text{total}}(\mathcal{G}) \text{ remains high}.
\]

The artifact becomes a terminal node of a larger optimization regime.
Its functionality persists,
but its grammar becomes inaccessible without authorization.

Compressive capture does not require coercion.
It arises from the coupling of efficiency incentives,
intellectual property regimes,
and networked infrastructure.

The result is a world in which objects function,
but do not explain themselves.

\subsection{The Instability of Externalized Grammar}

When interpretive keys are externalized, legibility becomes contingent
upon infrastructure continuity. Access to generative grammar depends not
on the artifact itself, but on the persistence of external systems that
mediate interpretation.

Under normal conditions this dependency remains largely invisible. The
artifact functions, diagnostic tools respond, and remote servers return
authorization. Under disruption, however, the consequences become
immediate: repair latency increases, substitution pathways narrow, and
local fabrication becomes infeasible. The object may remain physically
intact, yet become operationally inert.

This fragility is geometric rather than accidental. It arises from
collapsing the tolerance manifold without preserving a compensatory
interpretive channel within the artifact’s own structure. When slack is
removed and grammar is displaced, continuity becomes dependent on
external coordination rather than local recoverability.

\subsection{Toward Embedded Interpretation}

If compressive capture externalizes grammar,
the inverse strategy must embed it.

The task is not to eliminate optimization,
but to preserve sufficient slack within
the tolerance manifold to encode template-selection logic.

This motivates the introduction of
hyperpleonastic legibility:
the embedding of interpretive authority into geometry itself.

\section{The Geometry of Trust: The Spherical Selector}

In conventional industrial design, information is external to matter. 
A serial number is etched, a QR code is printed, a database entry is 
referenced, or a firmware key is embedded in silicon. The artifact 
itself does not explain its own composition or repair logic; it merely 
points outward. If the database disappears or the key is withheld, the 
object becomes epistemically mute. Its generative grammar is no longer 
recoverable from its physical presence.

This separation between artifact and interpretation constitutes a form 
of compressive capture. The object exists only as a dependent terminal 
of an external authority. Its dimensional legibility is artificially 
restricted.

To construct artifacts compatible with long-term civilizational 
alignment, we require a different principle: hyperpleonastic legibility. 
In such systems, the rule for interpreting the object is embedded within 
the object’s own geometry.

\subsection{The Canonical Example: The Quad-Ring Sphere}

Consider a sphere circumscribed by four colored rings. These rings are 
not decorative. They function as parametric selectors encoded directly 
into the symmetry of the form.

Let the inclination of each ring relative to a defined equatorial plane 
be denoted by
\[
\theta_1, \theta_2, \theta_3, \theta_4.
\]
Assume each $\theta_i$ is discretized into a finite set of admissible 
angular bins. The ordered quadruple
\[
(\theta_1, \theta_2, \theta_3, \theta_4)
\]
constitutes a physical coordinate in a public material-template grammar.

A particular configuration—for example
\[
(15^\circ, 30^\circ, 45^\circ, 90^\circ)
\]
—indexes a specific generative template describing the artifact’s
layer composition, assembly order, repair protocol, and recycling
conditions. The angular tuple functions as a coordinate within a
public template library, narrowing the space of possible grammars
to a determinate model class.

The sphere therefore does not directly encode full instructions.
Instead, it encodes a selector field that determines how the object
is to be interpreted. The information embedded in its geometry is not
the complete manual but the rule for retrieving the correct manual.
Interpretation remains structured rather than arbitrary, yet the
structure is materially anchored in the artifact itself.

\subsection{Holographic Recoverability}

Because each ring wraps the entire surface, the selector field is 
geometrically redundant. Even if the sphere is fractured or partially 
eroded, surviving arcs suffice to reconstruct the original angles. 
The decoding key is distributed across the object’s topology.

Destruction of legibility requires destruction of the artifact itself. 
The manual cannot be removed independently of the matter.

\subsection{Alignment Implications}

This architecture produces several structural consequences.

First, it increases dimensional control without increasing centralized 
bandwidth. No network query is required to classify the object. The 
artifact advertises its own interpretive context.

Second, it supports layered interpretability. At minimal perceptual 
resolution, the colors provide coarse classification. At higher 
resolution, angular measurement yields precise template identification. 
At still higher resolution, pigment microstructure may encode batch 
lineage or manufacturing history. Expertise scales continuously with 
observational refinement.

Third, it enforces civilizational incompressibility. The generative 
grammar remains reconstructible from the physical substrate. As long 
as the object exists, the interpretive pathway remains available. 
Information is not centralized in a removable database but distributed 
within geometry itself.

\subsection{Design Principle}

We may therefore state a preliminary alignment principle:

\medskip

\noindent
\textbf{Spherical Selector Principle.}
An artifact is repair-aligned if its template-selection logic is 
recoverable from macroscopic geometric features within human perceptual 
bandwidth, without requiring proprietary authorization.

\medskip

This principle reframes robustness. Durability is not merely resistance 
to fracture. It is persistence of legibility. An aligned artifact is one 
that continues to explain itself as long as it continues to exist.

\section{Selector Fields and the Tolerance Manifold}

The spherical example illustrates a broader structural principle. 
The essential feature is not the sphere, nor the colored rings, but 
the embedding of a template-selection logic within admissible 
geometric variation.

We now generalize this architecture.

\subsection{Functional Envelope and Tolerance Manifold}

Let an artifact be defined by a functional envelope $\mathcal{F}$ --- 
the set of physical configurations within which the object satisfies 
its mechanical, thermal, or structural requirements.

Within $\mathcal{F}$ there exists a subset of permissible micro-variations 
that do not degrade function. We denote this admissible variation space 
as the \emph{tolerance manifold} $\mathcal{T} \subset \mathcal{F}$.

In conventional engineering practice, $\mathcal{T}$ is treated as noise. 
Variation within $\mathcal{T}$ is minimized or ignored, as long as 
functional constraints are preserved.

In a hyperpleonastic system, $\mathcal{T}$ is reinterpreted as latent 
channel capacity.

\subsection{Definition of a Selector Field}

A \emph{Selector Field} is a structured mapping

\[
S : \mathcal{T} \rightarrow \mathcal{I}
\]

where $\mathcal{I}$ is a discrete index set corresponding to a public 
generative grammar of templates.

The selector field does not encode full instructions. It encodes 
which interpretive template must be applied.

An artifact is thus described by the pair
\[
(\mathcal{F}, S).
\]

The functional envelope guarantees operation. 
The selector field guarantees legibility.

\subsection{Recoverability Constraint}

For a selector field to satisfy repair-alignment criteria, it must meet 
a recoverability condition.

Let $\mathcal{O}$ denote the observational space available to a user 
within ordinary perceptual or low-instrument bandwidth. Let 
\[
R : \mathcal{O} \rightarrow \mathcal{T}
\]
be the reconstruction map from observable features to the tolerance 
manifold.

We require that the composite map
\[
S \circ R : \mathcal{O} \rightarrow \mathcal{I}
\]
be surjective onto the relevant subset of template indices.

In plain terms: the template-selection logic must be inferable from 
observable geometry without proprietary keys.

\subsection{Layered Interpretability}

A robust selector field supports multiple reconstruction resolutions.
Interpretive access need not be uniform; it may unfold across increasing
levels of observational precision.

Let $\mathcal{O}_1 \subset \mathcal{O}_2 \subset \mathcal{O}_3$ denote
nested observational resolutions corresponding, for example, to visual
inspection, measurement with simple tools, and microscopic or spectral
analysis. Each level enlarges the recoverable portion of the generative
grammar while remaining consistent with lower-resolution readings.

At the coarsest level, the selector field may permit categorical
classification. At intermediate resolution, it may permit parametric
specification. At the finest level, it may permit batch-level lineage or
material calibration. The key structural property is monotonic
refinement: higher-resolution observation never contradicts lower-level
interpretation but instead extends it.

Layered interpretability therefore ensures that legibility is not
binary. An artifact does not transition from opaque to transparent in a
single step; rather, it reveals progressively finer structure as
observational capacity increases.

We require monotonic refinement:
\[
S \circ R_1 \subseteq S \circ R_2 \subseteq S \circ R_3.
\]

Higher resolution may increase specificity but must not contradict 
lower-resolution interpretation.

This condition ensures that expertise refines meaning rather than 
overturning it. The novice classification remains valid even when 
expert detail is added.

\subsection{Slack as Channel Capacity}

Let $\dim(\mathcal{T})$ denote the dimensionality of the tolerance 
manifold. This dimensionality represents degrees of freedom that do 
not compromise functionality.

If $\mathcal{T}$ is collapsed to a single canonical configuration, 
its channel capacity approaches zero. The artifact becomes silent.

If instead variation within $\mathcal{T}$ is discretized into bins, 
the cardinality of $\mathcal{I}$ is bounded above by the effective 
quantization of $\mathcal{T}$ under functional constraints.

Thus, robustness implies potential encoding capacity.

We may state the general principle:

\medskip

\noindent
\textbf{Selector Field Theorem (Preliminary Form).}
Any artifact possessing a nontrivial tolerance manifold 
$\mathcal{T}$ admits a selector field mapping into a finite template 
grammar without degrading functionality, provided encoding amplitude 
remains within admissible bounds.

\medskip

The theorem does not require exotic materials. It follows directly 
from the existence of bounded slack in engineered systems.

\subsection{Alignment Consequence}

Conventional optimization seeks to minimize $\dim(\mathcal{T})$ in the 
name of efficiency, aesthetic uniformity, or cost control. This reduces 
encoding capacity and increases dependence on external documentation.

Hyperpleonastic design preserves $\dim(\mathcal{T})$ and assigns it 
semantic function.

Alignment, in this context, is the deliberate refusal to collapse 
the tolerance manifold to zero informational depth.

\section{Selector Fields and Civilizational Geometry}

The preceding analysis treats selector fields as a property of individual 
artifacts. We now situate this mechanism within the larger problem of 
civilizational alignment under asymmetric optimization.

\subsection{Dimensional Control Surfaces}

Let $D_{\text{control}}$ denote the effective dimensional control surface
available to human agents interacting with a system. This quantity
measures the number of independent parameters through which users can
classify the artifact, diagnose its state, intervene in its function,
and restore or modify its configuration. It represents not merely
access, but structured agency: the degree to which meaningful variation
is accessible for inspection and adjustment.

In opaque systems, $D_{\text{control}}$ is artificially restricted.
Critical parameters are concealed behind proprietary firmware, sealed
enclosures, or centralized documentation servers. Although the artifact
exists physically, its generative grammar is displaced into external
infrastructure. The object functions as a terminal interface rather than
as a legible substrate.

Selector fields expand $D_{\text{control}}$ without increasing
centralized coordination bandwidth. By embedding interpretive structure
within the artifact’s own geometry, they restore local recoverability of
the key. Agency is increased not by adding complexity, but by preserving
dimensional access within admissible physical variation.

\subsection{Alignment Under Temporal Asymmetry}

In earlier work, we described a Nyquist-style containment constraint of the form

\[
\left( \frac{C_{\text{swarm}}}{D_{\text{control}}} \right)
\left( \frac{T_{\text{human}}}{\tau_{\text{AI}}} \right) \leq 1.
\]

This inequality expresses a structural requirement: as system capability 
$C_{\text{swarm}}$ increases and machine action cycles $\tau_{\text{AI}}$ 
accelerate, the available dimensional control surface must increase 
proportionally to prevent endogenous governance.

Selector fields act directly on $D_{\text{control}}$.

They do not slow $\tau_{\text{AI}}$.
They do not reduce $C_{\text{swarm}}$.
They increase interpretive dimensionality at the level of the artifact.

In doing so, they increase damping capacity within the inequality.

\subsection{Opacity as Compression}

When tolerance manifolds are collapsed to canonical uniformity, 
interpretive dimensionality decreases. The object becomes dependent on 
external keys. This compression reduces $D_{\text{control}}$ while leaving 
optimization capacity intact.

Such systems are dynamically fragile. Repair becomes centralized. 
Modification becomes licensed. Knowledge becomes revocable.

Compressive capture is therefore not merely informational; it is geometric.

\subsection{Distributed Damping}

A civilization composed of selector-field artifacts exhibits distributed
damping characteristics. Stability does not depend exclusively on
centralized intervention; instead, it arises from the capacity of local
agents to interpret and act upon the generative grammars embedded within
their tools.

In such a system, agents can reconstruct template logic, diagnose
failure modes, restore functionality, and replicate components without
requiring high-bandwidth coordination with a central authority. Each
artifact carries sufficient interpretive structure to permit bounded
autonomy in repair and reconstruction.

Damping, in this framework, is not imposed from above as a regulatory
overlay. It emerges from distributed legibility. The capacity to absorb
perturbation is built into the geometry of artifacts themselves, rather
than enforced through remote supervision or institutional latency.

This reduces phase lag in oversight loops. Intervention does not require 
external authentication or proprietary mediation. Governance remains 
exogenous to automated subsystems rather than collapsing into them.

\subsection{Incompressibility as Persistence of Legibility}

We may now refine the notion of civilizational incompressibility.

A civilization is compressible when its generative grammars are 
externalized into revocable databases, sealed firmware, or proprietary 
cloud dependencies.

It is incompressible when the rules required for repair, replication, 
and reinterpretation remain embedded within physical or publicly 
distributed structures.

Selector fields operationalize this embedding.

They ensure that artifacts do not become mute terminals of external 
optimization regimes.

\subsection{Reframing Robustness}

Within this framework, robustness admits two distinct meanings.
Structural robustness refers to resistance against physical degradation,
the capacity of an artifact to maintain functional integrity under
stress, wear, or environmental exposure. Interpretive robustness, by
contrast, refers to the persistence of legibility under degradation—the
capacity of an artifact to continue revealing its generative grammar even
when damaged.

Selector fields primarily enhance interpretive robustness. Their role is
not merely to preserve structural strength but to preserve recoverable
meaning. Even when partially damaged, the artifact continues to advertise
its template logic. As long as sufficient geometric residue remains, the
generative grammar remains inferable from the substrate itself.

The object does not merely endure; it continues to explain itself.

\subsection{Alignment as Geometric Constraint}

We may therefore restate alignment at the material level:

\medskip

\noindent
\textbf{Geometric Alignment Principle.}
An artifact is alignment-compatible if its generative grammar remains 
reconstructible from its physical substrate under bounded degradation, 
without reliance on proprietary authorization.

\medskip

This principle does not prohibit advanced optimization. It constrains 
how optimization may externalize interpretive authority.

The refusal to collapse tolerance manifolds is thus not aesthetic excess. 
It is a geometric safeguard against dimensional asymmetry.

\section{Thermovariable Membranes and Wearable Control Surfaces}

Wearable thermal regulation systems occupy a high-stakes domain. 
They participate directly in human homeostasis. A garment that widens 
its weave in response to heat or actively circulates warmed fluid 
is not merely an accessory; it is part of a distributed physiological loop.

If such systems are opaque, the wearer possesses no effective 
dimensional control over their own thermal regulation. Failure 
modes become silent. Diagnosis becomes proprietary. Repair becomes 
externalized.

We now examine how selector-field design applies to thermovariable 
membranes and electronic textiles.

\subsection{Bimetallic Weaves as Mechanical Logic}

Consider a textile incorporating bimetallic strips woven into its 
structure. A bimetallic element bends predictably as temperature 
changes due to differential expansion coefficients.

In conventional ``smart'' garments, these elements are laminated 
within sealed layers. Their state is invisible without instrumentation.

In a hyperpleonastic membrane, the bending action itself becomes 
part of the selector field.

Let the curvature state of a strip at temperature $T$ be denoted 
by $\kappa(T)$. The textile weave is arranged such that changes in 
$\kappa(T)$ reveal color-coded underlayers.

At temperature $T_1$, the weave remains in a closed configuration,
exposing substrate $C_1$ (for example, a blue underlayer). At a higher
temperature $T_2 > T_1$, thermally induced bending widens the apertures
of the weave, revealing substrate $C_2$ (for example, a red underlayer).
The chromatic transition therefore tracks the mechanical state of the
membrane as a direct function of temperature.

The operational state of the membrane is consequently geometrically
self-evident. No software query is required, no remote synchronization
is necessary, and no proprietary diagnostic interface mediates
interpretation. The control surface is perceptually available, embedded
in the visible transformation of the material itself.

\subsection{Embedded Thermal and Fluidic Garments}

Active heating or cooling garments may incorporate resistive heating
traces, fluid-circulating microtubing, and voltage regulation nodes
within their structure. These components form a distributed thermal
control system that operates in close contact with the human body.

In opaque implementations, such subsystems are concealed within sealed
laminates or hidden between fabric layers. Electrical pathways and
fluid channels become inaccessible to inspection. When failure occurs,
it presents as binary dysfunction: the garment either heats or cools,
or it does not, with no visible indication of the underlying fault.

A hyperpleonastic design instead arranges these subsystems along
visible braiding paths that encode operational constraints. The routing
of conductive traces and fluid channels is not merely functional but
interpretive. Their geometry advertises voltage limits, pressure bounds,
and safe modification pathways, transforming the garment from a sealed
device into a legible thermal instrument.

Let the braid pitch $p$ encode maximum voltage $V_{\max}$, and 
let the spacing ratio $s$ encode maximum fluid pressure $P_{\max}$.

These parameters remain within the tolerance manifold of safe 
mechanical performance. They do not degrade functionality. 
They serve as selector-field coordinates.

If a break occurs, the garment advertises both its type and its 
operational limits through visible structure.

Diagnosis migrates from proprietary diagnostic firmware to 
direct geometric inspection.

\subsection{Homeostatic Sovereignty}

Wearable systems form a feedback loop:

\[
\text{Body} \leftrightarrow \text{Garment} \leftrightarrow \text{Environment}
\]

If the garment's generative grammar is externalized into proprietary 
software or remote optimization services, the wearer’s effective 
control dimension collapses.

Selector-field garments maintain homeostatic sovereignty by 
embedding interpretive logic into the substrate.

The user need not trust remote infrastructure to understand 
whether the membrane is open, closed, shorted, or obstructed.

\subsection{Failure Modes and Interpretive Robustness}

Consider two garments possessing identical functional envelopes and
comparable thermal performance under ideal conditions. In the first
case, heating traces are concealed within sealed layers, and electrical
failures produce no externally legible indication of fault. In the
second, the heating traces are arranged along visible braiding paths,
with braid geometry encoding voltage class and operational limits.

At the outset, both garments may perform equivalently. Their structural
robustness and thermal capacity may be indistinguishable under nominal
operation. Under degradation, however, interpretive robustness diverges.

In the opaque case, failure manifests as silence. The garment ceases to
function, yet provides no structural indication of the underlying
disturbance. In the hyperpleonastic case, degradation remains legible.
Distortions in braid geometry, discontinuities in conductive pathways,
or chromatic anomalies reveal the location and character of failure.
The artifact does not merely stop; it continues to communicate.

The divergence lies not in mechanical strength but in the persistence
of recoverable grammar under stress. One garment becomes mute; the
other continues to explain itself.

This distinction maps directly onto dimensional control:

\[
D_{\text{control}}^{\text{opaque}} < D_{\text{control}}^{\text{selector}}.
\]

The inequality persists even under partial failure.

\subsection{The Black-Box Safety Tax}

Systems that collapse their tolerance manifolds and conceal their 
selector fields impose a systemic liability.

Because diagnosis requires centralized expertise, response latency 
increases. Small failures escalate into catastrophic events.

Insurance, infrastructure reliability, and public safety models 
implicitly penalize opacity.

Selector-field garments reduce this opacity tax by embedding 
legibility within the artifact.

\subsection{Didactic Surfaces}

Hyperpleonastic membranes function not only as regulatory systems but
as didactic environments. Their interpretive structure unfolds across
levels of expertise. A novice wearer perceives chromatic transitions
indicating thermal state; a technician measures braid pitch to infer
voltage class and operational limits; a specialist traces conductive
paths to reconstruct circuit topology and diagnose localized faults.
Each level of observation yields coherent information appropriate to
its resolution.

The garment thereby becomes both instrument and instructor. It does not
merely perform a function; it communicates the logic of that function
through its geometry. Interpretation is not mediated exclusively by
external documentation but emerges from structured interaction with the
material substrate.

This property reinforces civilizational incompressibility. Knowledge
remains coupled to matter rather than sequestered in archival systems or
proprietary databases. As long as the artifact persists, the grammar of
its operation remains partially recoverable from its form.

\section{Crisis Geometry: Selector Fields Under Infrastructure Failure}

The practical significance of selector-field design becomes visible 
under conditions of infrastructural stress. We therefore examine 
a simplified crisis scenario and compare system behavior in opaque 
and hyperpleonastic regimes.
\subsection{Scenario: Heatwave and Network Failure}

Assume a prolonged heatwave coincides with intermittent electrical and
network outages. Citizens rely on thermovariable garments for
temperature regulation, placing thermal control within a constrained and
potentially fragile infrastructure loop.

Two distinct design architectures coexist in this environment. In the
first, thermal control is mediated by embedded firmware and
cloud-synchronized optimization routines. Regulation depends upon remote
coordination and periodic network exchange. In the second, thermal
response is governed by substrate-defined mechanical logic and visible
geometric encoding. Regulation emerges directly from material structure,
without dependence on external synchronization.

Let the average machine action cycle be $\tau_{\text{AI}}$, and let human
diagnosis time be $T_{\text{human}}$. When network infrastructure is
stable, firmware-mediated garments may function normally, with control
loops operating at machine timescales. Under network failure, however,
remote control loops are severed. The garment’s regulatory logic no
longer receives external updates, and its operational state may become
ambiguous or frozen.

By contrast, in substrate-governed systems, thermal regulation remains
locally coupled to physical conditions. Mechanical response proceeds
without reference to remote computation, and interpretive access remains
available through direct inspection. The divergence is not merely
technical but geometric: one architecture collapses when coordination
channels fail, whereas the other retains autonomy within its material
envelope.

\subsection{Control Surface Collapse}

For opaque garments, effective dimensional control reduces to:

\[
D_{\text{control}}^{\text{opaque}} \rightarrow 0
\]

once firmware locks or cloud dependencies prevent local override.

Diagnosis requires specialized equipment or manufacturer access. 
Human intervention time $T_{\text{human}}$ increases dramatically.

Under crisis conditions, the Nyquist-style constraint degrades:

\[
\left( \frac{C_{\text{system}}}{D_{\text{control}}^{\text{opaque}}} \right)
\left( \frac{T_{\text{human}}}{\tau_{\text{failure}}} \right) \gg 1.
\]

Oversight becomes endogenous to unavailable infrastructure.

\subsection{Selector-Field Persistence}

For selector-field garments, the thermal logic remains substrate-defined.

Mechanical widening of weave occurs independently of network state.
Voltage class and pressure limits remain encoded in braid geometry.
Failure modes remain visible.

Thus,

\[
D_{\text{control}}^{\text{selector}} > 0
\]

even when external coordination fails.

Human intervention time remains bounded:

\[
T_{\text{human}}^{\text{selector}} \approx T_{\text{inspection}} + T_{\text{manual repair}}.
\]

The control inequality degrades far less severely.

\subsection{Graceful Degradation}

We define \emph{graceful degradation} as the persistence of interpretable 
structure under partial system failure.

Opaque systems degrade discretely: functional $\rightarrow$ nonfunctional.

Selector-field systems degrade continuously: optimal $\rightarrow$ reduced efficiency $\rightarrow$ repairable state.

This continuity arises because the selector field distributes interpretive 
authority into geometry rather than central firmware.

\subsection{Supply Chain Disruption}

Consider also a supply-chain interruption preventing access to proprietary 
replacement components.

In opaque regimes, part identification requires SKU databases or embedded 
digital authentication.

In selector-field regimes, the template index remains recoverable from 
visible geometry. Local fabrication becomes feasible.

Let $L$ denote localization capacity.

Then under disruption:

\[
L^{\text{opaque}} \rightarrow 0, \qquad
L^{\text{selector}} > 0.
\]

Localization capacity functions as civilizational damping.

\subsection{Thermodynamic Stability and Interpretive Entropy}

Under crisis, entropy increases: infrastructure fails, coordination 
fragments, resources become scarce.

If interpretive grammars are externalized, informational entropy 
increases faster than physical entropy. Objects become unintelligible 
before they become unusable.

Selector-field systems resist this interpretive entropy.

As long as geometric residues remain, template selection remains possible.

Legibility decays more slowly than infrastructure.

\subsection{Crisis Conclusion}

The selector field is therefore not a luxury feature. 
It is a stabilizing geometric constraint.

Under normal conditions, it increases agency.
Under crisis conditions, it preserves agency.

Civilizational incompressibility is revealed not as ideological 
abstraction, but as the persistence of local interpretability 
when centralized optimization loops fail.

\section{Recoverable Information Density}

Selector-field artifacts preserve legibility under degradation. 
We now formalize this property as Recoverable Information Density (RID).

\subsection{Total Generative Information}

Let an artifact be produced according to a generative grammar
$\mathcal{G}$. The total information required to fully specify its
construction, repair logic, and material composition is

\[
I_{\text{total}} = H(\mathcal{G}),
\]

where $H(\cdot)$ denotes the Shannon information required to describe
the grammar.

This quantity encompasses the specification of material ratios,
assembly order, stress tolerances, repair constraints, and operational
limits. It represents the complete informational description necessary
to reproduce or restore the artifact under ideal interpretive
conditions.

In most contemporary systems, however, $I_{\text{total}}$ resides
primarily in external documentation, firmware repositories, or
proprietary databases. The artifact itself carries only a minimal
projection of its generative grammar, insufficient for full
reconstruction without access to external informational infrastructure.

\subsection{Recoverable Information}

Let $\mathcal{O}$ denote the observational space available to a user 
under bounded instrumentation (e.g., unaided perception, simple tools).

Let $I_{\text{recoverable}}$ be the mutual information between 
$\mathcal{O}$ and the generative grammar $\mathcal{G}$:

\[
I_{\text{recoverable}} = I(\mathcal{O}; \mathcal{G}).
\]

This quantity measures how much of the artifact's generative grammar 
can be inferred from its physical substrate without proprietary keys.

\subsection{Definition of RID}

We define Recoverable Information Density as

\[
\text{RID} = \frac{I_{\text{recoverable}}}{I_{\text{total}}},
\]

so that

\[
0 \leq \text{RID} \leq 1.
\]

The limiting case $\text{RID} = 0$ corresponds to a fully opaque
artifact, one whose generative grammar is not materially recoverable
from inspection of the artifact itself. The opposite limit,
$\text{RID} = 1$, corresponds to a fully self-describing artifact in
which the entirety of its generative grammar is inferable from its
physical substrate without reliance on external keys.

In practice, most systems occupy an intermediate regime in which a
fraction of the generative grammar remains locally recoverable while
the remainder is externalized. The value of $\text{RID}$ therefore
quantifies the degree to which interpretive authority is embedded within
the artifact rather than displaced into surrounding infrastructure.

\subsection{RID and Dimensional Control}

We may relate RID to dimensional control surface:

\[
D_{\text{control}} \propto I_{\text{recoverable}}.
\]

As RID increases, the effective number of independent parameters 
accessible to the user increases.

Selector fields increase RID by mapping tolerance manifold variation 
onto template indices.

\subsection{The Black-Box Tax}

Let $\Phi$ denote the opacity penalty associated with an artifact. 
We define:

\[
\Phi = 1 - \text{RID}.
\]

High $\Phi$ corresponds to high dependence on centralized authority 
for interpretation.

Under infrastructure failure, artifacts with high $\Phi$ exhibit
increased repair latency, elevated probability of catastrophic failure,
and reduced localization capacity. Because interpretive keys are
externalized, restoration depends on the reestablishment of coordination
channels rather than on local inspection and intervention.

Opacity therefore imposes systemic risk externalities. The cost of
interpretive displacement is not confined to the individual artifact
but propagates through the surrounding network of dependencies,
amplifying fragility under disruption.

\subsection{Degradation and RID Persistence}

Let $d$ denote degradation level (e.g., surface erosion, partial damage).

Define $\text{RID}(d)$ as the recoverable information density 
under degradation.

Selector-field artifacts are designed to satisfy:

\[
\frac{d}{dd} \text{RID}(d) \approx 0 \quad \text{for small } d.
\]

That is, legibility decays slowly relative to structural degradation.

Opaque artifacts satisfy:

\[
\frac{d}{dd} \text{RID}(d) \ll 0,
\]

indicating rapid collapse of interpretive capacity under minor damage.

\subsection{RID as a Regulatory Metric}

Recoverable Information Density provides a measurable criterion 
for policy and design evaluation.

Regulatory standards could require:

\[
\text{RID} \geq \rho_{\text{min}}
\]

for critical infrastructure artifacts.

Such thresholds would not mandate disclosure of proprietary 
templates, but would require that template-selection logic 
remain physically inferable.

\section{An Alignment Stability Theorem}

We now unify the preceding constructions into a single stability claim.
The aim is not to assert a universal law of governance, but to formalize
a structural mechanism by which material legibility enlarges human
control surfaces and thereby increases the stability margin of oversight
under asymmetric optimization.

\subsection{Setup}

Let an artifact be produced by a generative grammar $\mathcal{G}$ with
total description complexity
\[
I_{\text{total}} := H(\mathcal{G}).
\]

Let $\mathcal{T}$ be the tolerance manifold of admissible variation
within the functional envelope, and let $S : \mathcal{T} \to \mathcal{I}$
be a selector field into a discrete template index set $\mathcal{I}$.

Let $\mathcal{O}$ be the bounded observational space available to a user
without proprietary authorization, and let $R:\mathcal{O}\to\mathcal{T}$
be the reconstruction map. The observable template index is then
\[
\widehat{\imath} := (S\circ R)(o)\in \mathcal{I}.
\]

Define recoverable information density
\[
\text{RID} := \frac{I(\mathcal{O};\mathcal{G})}{H(\mathcal{G})}
= \frac{I_{\text{recoverable}}}{I_{\text{total}}}.
\]

Finally, let $D_{\text{control}}$ denote the effective dimensional control
surface available to humans for diagnosis and intervention in the relevant
sociotechnical system. We assume that, over the regime of interest,
\[
D_{\text{control}} \ \asymp\  D_0 + \lambda\, I_{\text{recoverable}},
\]
where $D_0$ is baseline actuation capacity (tools, skills, institutions)
and $\lambda>0$ converts recoverable information into usable degrees of
intervention.

\subsection{A Containment Constraint}

We adopt the Nyquist-style containment constraint previously introduced:
\[
\left( \frac{C_{\text{system}}}{D_{\text{control}}} \right)
\left( \frac{T_{\text{human}}}{\tau_{\text{system}}} \right)\ \leq\ 1,
\]
where $C_{\text{system}}$ is aggregate system capability, $T_{\text{human}}$
is the human decision cycle time for meaningful oversight, and
$\tau_{\text{system}}$ is the system action cycle time (including automated
subsystems).

The left-hand side expresses a dimensionless instability pressure:
capability per controllable dimension, multiplied by temporal mismatch.

\subsection{Theorem: Legibility Increases the Stability Margin}

\medskip
\noindent
\textbf{Theorem (Alignment Stability via Recoverable Legibility).}
Assume a class of artifacts in critical loops admits nontrivial tolerance
manifolds $\mathcal{T}$ and selector fields $S$ such that, for bounded
instrumentation $\mathcal{O}$, the observable template index
$\widehat{\imath}=(S\circ R)(o)$ is recoverable with low ambiguity and
$\text{RID} \geq \rho$ for some $\rho>0$.
Then, holding $C_{\text{system}}$, $T_{\text{human}}$, and $\tau_{\text{system}}$
fixed, the containment inequality admits a strictly larger stability margin
than in the opaque regime $\text{RID}\approx 0$.

Equivalently, for fixed temporal mismatch, the maximum admissible capability
before instability scales as
\[
C_{\text{system}}^{\max}\ \propto\ D_{\text{control}}
\ \asymp\ D_0 + \lambda\, \rho\, I_{\text{total}}.
\]

\medskip

\noindent
\emph{Proof (structural).}
By definition, $\text{RID}\ge \rho$ implies
\[
I_{\text{recoverable}} = I(\mathcal{O};\mathcal{G}) \ge \rho\, I_{\text{total}}.
\]
Under the monotone mapping assumption
$D_{\text{control}} \asymp D_0 + \lambda I_{\text{recoverable}}$,
we obtain
\[
D_{\text{control}} \ge D_0 + \lambda\, \rho\, I_{\text{total}}.
\]
Substituting into the containment inequality yields a smaller (or equal)
left-hand side relative to the opaque case in which $\rho\approx 0$ and
$D_{\text{control}}\approx D_0$. Therefore the stability margin increases.
\hfill $\square$

\subsection{Corollaries}

\paragraph{Corollary 1 (Opacity as a Stability Liability).}
Define the black-box tax $\Phi:=1-\text{RID}$. For fixed $I_{\text{total}}$,
an increase in $\Phi$ decreases $D_{\text{control}}$ and therefore lowers
$C_{\text{system}}^{\max}$ before instability. Infrastructures dominated by
high-$\Phi$ artifacts are dynamically fragile under optimization pressure.

\paragraph{Corollary 2 (Crisis Persistence).}
If selector-field artifacts satisfy a slow-decay condition
$\text{RID}(d)\ge \rho$ for degradation levels $d$ in a crisis regime,
then $D_{\text{control}}$ remains bounded away from collapse even as
physical damage accumulates. This yields graceful degradation rather than
binary failure.

\paragraph{Corollary 3 (Localization Capacity).}
If template indices $\widehat{\imath}$ are recoverable locally, then the
feasible set of local interventions expands to include distributed
fabrication and substitution. Localization capacity is therefore an
increasing function of RID.

\subsection{Interpretation}

The theorem does not claim that legibility solves alignment in general.
It specifies a mechanism: selector fields convert slack in the tolerance
manifold into recoverable interpretive structure, thereby increasing
human control dimensionality. This enlarges the stability margin of
oversight constraints under temporal mismatch.

On this view, ``alignment'' is not only a property of policies and
optimizers. It is a property of the material substrate on which
governance must act.

\subsection{Design and Policy Implication}

A direct consequence is that increasing $D_{\text{control}}$ need not
require increased surveillance, higher-bandwidth bureaucracy, or tighter
central command. It can be achieved by embedding interpretive authority
into the artifact itself.

In regulatory terms, minimum-threshold requirements of the form
\[
\text{RID}\ \ge\ \rho_{\min}
\]
become stability constraints rather than consumer conveniences.
They preserve exogenous oversight by preventing the systematic collapse
of control dimensionality via opacity.

\section{Domestic Appliances and Everyday Incompressibility}

The selector-field framework is not limited to high-performance
materials or wearable systems. Its relevance is clearest in the
domain of ordinary household infrastructure.

Milk production, yogurt fermentation, and paper fabrication are
historically domestic activities. In industrial societies, they
have been externalized into centralized supply chains. The household
becomes a terminal node consuming opaque outputs.

We examine how selector-field design alters this dynamic.

\subsection{Milk and Yogurt Systems}

Consider a domestic yogurt appliance.

In its opaque form, the device contains a sealed heating element,
a thermostat governed by proprietary firmware, a hidden temperature
sensor, and a non-serviceable power module. These components are
physically present but structurally inaccessible, and their operational
relationships are not directly inferable from inspection.

Failure presents as binary. The device ceases to function without
providing legible indication of the underlying fault. The fermentation
grammar, including temperature curves, inoculation timing, and thermal
inertia, is externalized into manuals or embedded firmware rather than
remaining recoverable from the physical substrate.

Recoverable Information Density is low:
\[
\text{RID}_{\text{opaque}} \approx 0.
\]

\subsubsection{Selector-Field Fermentation}

In a hyperpleonastic configuration, the fermentation vessel integrates
interpretive structure directly into its physical form. A visible
bimetallic indicator strip is calibrated to the relevant fermentation
temperature ranges, a color-coded thermal ring marks incubation phase,
a mechanical timer exposes its gearing ratio in a manner that encodes
fermentation class, and heating traces are routed visibly with clear
voltage class markings. Each structural element performs a functional
role while simultaneously participating in a selector field.

The device therefore advertises its target temperature window, phase
transitions, power limits, and repair access points through geometric
and chromatic cues rather than through external documentation. The
grammar of fermentation is not displaced into firmware or manuals but
remains partially recoverable from inspection of the vessel itself.

Even in the absence of formal documentation, the generative grammar of
the fermentation process remains inferable. Formally,

\[
\text{RID}_{\text{selector}} > 0.
\]

The distinction is not one of convenience but of interpretive
persistence. The device continues to communicate its operational logic
under conditions in which external informational infrastructure may be
absent.

\subsection{Toilet Paper and Material Provenance}

Toilet paper appears trivial under ordinary conditions. During supply
disruptions, however, it becomes a focal point of systemic fragility,
revealing the dependence of everyday materials on centralized
production and distribution networks.

A hyperpleonastic paper product could encode material provenance
directly into its fiber orientation patterning and watermark geometry.
Variations in weave direction, density modulation, and watermark
symmetry could communicate fiber composition ratios, source pulp class,
biodegradation rate, and recommended alternative fiber substitutions.
These signals need not disclose proprietary sourcing contracts; they
need only expose sufficient material grammar to enable replication,
adaptation, or substitution under constrained conditions.

In such a system, local paper production becomes reconstructible from
the substrate itself. The artifact does not merely serve its immediate
use but participates in a distributed template library for material
regeneration. Localization capacity correspondingly increases, not by
mandate but by embedded legibility.

\subsection{From Consumer to Maintainer}

In opaque domestic infrastructure, the household functions as a passive
endpoint within a larger optimization network. Appliances operate, but
their generative grammars remain externalized. Diagnosis depends on
authorized service channels, repair requires proprietary components,
and substitution pathways are constrained by centralized supply chains.
Agency is limited to consumption and replacement.

In selector-field domestic infrastructure, this condition changes
structurally. The household becomes diagnostician, repairer, replicator,
and substitutor. Interpretive cues embedded within artifacts permit
local identification of failure modes, reconstruction of template logic,
and adaptation using available materials. The transition is not merely
cultural but geometric; the control surface expands because grammar is
embedded rather than displaced.

Formally, this shift corresponds to an increase in
$D_{\text{control}}$ at the smallest scale of civil society. Agency is
not granted through policy alone but materialized through the
recoverability of generative structure within the objects of everyday
life.

\subsection{Dimensional Poverty in the Home}

When appliances externalize their generative grammars into firmware
and proprietary supply chains, the household experiences
Dimensional Poverty.

Formally:

\[
D_{\text{control}}^{\text{household}} \rightarrow D_0,
\]

where $D_0$ represents minimal actuation without interpretive access.

Selector-field domestic tools increase this baseline.

The home becomes a node of distributed damping rather than
a dependency sink.

\subsection{Everyday Alignment}

Alignment is often framed in terms of extreme scenarios:
superintelligence, autonomous weapons, catastrophic risk.

Yet civilizational stability depends equally on mundane resilience.

Milk fermentation and paper fabrication are low-energy,
high-frequency processes. When such processes remain locally
interpretable, the cumulative effect is substantial.

Civilizational incompressibility is not achieved through grand
infrastructure alone. It is achieved through the persistence of
legibility in everyday artifacts.

\section{Dimensional Poverty}

The preceding sections have defined Recoverable Information Density (RID)
at the artifact level and related it to the dimensional control surface
$D_{\text{control}}$. We now generalize this framework to the civic scale.

\subsection{Household Control Surface}

Let a household $H$ interact with a set of artifacts
\[
\mathcal{A} = \{A_1, A_2, \dots, A_n\}.
\]

Each artifact $A_i$ has recoverable information density $\text{RID}_i$
and contributes to the household’s effective control dimensionality.

Define household dimensional control as:

\[
D_H = \sum_{i=1}^{n} w_i \, I_{\text{recoverable}}(A_i),
\]

where $w_i$ weights the systemic importance of each artifact
(e.g., thermal systems weighted more heavily than decorative objects).

In opaque economies, $\text{RID}_i \approx 0$ for critical goods.
Thus,

\[
D_H \approx D_0,
\]

where $D_0$ represents baseline physical manipulation without
interpretive authority.

\subsection{Definition of Dimensional Poverty}

We define Dimensional Poverty as the condition in which:

\[
D_H < D_{\text{threshold}},
\]

where $D_{\text{threshold}}$ is the minimum dimensional control required
to sustain local repair, substitution, and autonomous operation under
moderate infrastructure disruption.

Dimensional Poverty is therefore not income poverty.
It is a contraction of interpretive agency.

A household may possess advanced devices yet remain
interpretively impoverished if those devices are opaque.

\subsection{Systemic Amplification}

Dimensional Poverty scales nonlinearly.

Let $\mathcal{H}$ denote a set of households in a region.
Define aggregate civic dimensionality:

\[
D_{\text{civic}} = \sum_{H \in \mathcal{H}} D_H.
\]

If artifacts across the region share identical opaque dependencies
(e.g., identical firmware locks, centralized cloud authentication),
then local failures propagate systemically.

Low $D_H$ across many households yields:

\[
D_{\text{civic}} \rightarrow 0 \quad \text{under shared failure modes}.
\]

Selector-field artifacts introduce heterogeneity and distributed recovery,
raising both $D_H$ and $D_{\text{civic}}$.

\subsection{Dimensional Inequality}

Just as income can be unevenly distributed, so can dimensional control.

Define dimensional inequality between agents $i$ and $j$ as:

\[
\Delta D_{ij} = |D_i - D_j|.
\]

In highly opaque systems, manufacturers and platform owners possess
high $D_{\text{control}}$, while users possess minimal $D_H$.

This asymmetry creates structural dependence independent of formal law.

Dimensional inequality is therefore a governance parameter,
not merely a design choice.

\subsection{Dimensional Poverty and the Black-Box Tax}

Recall the opacity penalty:

\[
\Phi = 1 - \text{RID}.
\]

We may define household black-box exposure as:

\[
\Phi_H = \sum_{i=1}^{n} w_i \Phi_i.
\]

High $\Phi_H$ implies that most generative grammars governing daily life
are externally controlled rather than materially embedded within
household artifacts. Interpretive authority resides in remote
documentation systems, proprietary firmware layers, or centralized
service networks rather than in the objects themselves.

Under crisis conditions, high $\Phi_H$ correlates with increased repair
latency, reduced substitution capacity, and elevated catastrophic
dependency. Because generative grammars are not locally recoverable,
restoration depends on reestablishing external coordination rather than
on distributed maintenance. Fragility therefore propagates not from
mechanical weakness alone but from the externalization of interpretive
structure.

\subsection{Alignment as a Civil Liberty}

If selector-field design increases $\text{RID}_i$ for essential artifacts
and thereby expands $D_H$, the aggregate dimensional control surface
available at the household level, then alignment can no longer be
understood exclusively as a problem of software safety or algorithmic
constraint. It becomes a structural property of material systems and a
condition of civic agency.

In this framework, alignment concerns the preservation of sufficient
civic dimensionality to prevent endogenous governance collapse. When
interpretive grammars are systematically externalized, the control
surface available to citizens contracts, and oversight becomes dependent
upon the very infrastructures it seeks to regulate. Stability then
depends on uninterrupted coordination rather than on distributed
legibility.

Dimensional Poverty may therefore be defined as the condition in which
agents lose the capacity to read, repair, and reproduce the material
substrate of daily life. It is not merely economic deprivation but
geometric contraction of interpretive space.

Civilizational incompressibility requires that $D_H$ remain bounded away
from zero for critical systems. As long as this lower bound is preserved,
generative grammars remain locally recoverable, and material governance
does not collapse into opaque dependency.

\section{Opaque and Hyperpleonastic Economies}

We now compare two idealized economic regimes distinguished not by
intention but by structure. In the first, which we term an opaque
economy, generative grammars are systematically externalized into
proprietary databases, sealed firmware layers, and centralized supply
chains. Artifacts function as terminal interfaces whose interpretive
keys reside elsewhere. Coordination, repair, and replication depend on
access to institutional repositories rather than on material inspection.

In the second regime, which we term a hyperpleonastic economy, selector
fields are embedded within artifacts themselves, preserving recoverable
information density across distributed nodes. Generative grammars are
not fully disclosed, but they are partially encoded in geometry, texture,
and admissible variation. Interpretation remains locally bounded yet
structurally available.

The distinction between these regimes is not moral but geometric. It
concerns the distribution of generative description length, the location
of interpretive authority, and the dimensionality of control surfaces
available to agents.

\subsection{Aggregate Dimensional Control}

Let $\mathcal{H}$ denote the set of households, and let 
$\mathcal{I}$ denote industrial producers.

Define aggregate dimensional control:

\[
D_{\text{total}} = D_{\text{civic}} + D_{\text{industrial}}.
\]

In opaque economies:

\[
D_{\text{civic}} \ll D_{\text{industrial}}.
\]

Control dimensionality is concentrated at production nodes.
Consumers possess minimal interpretive authority.

In hyperpleonastic economies:

\[
D_{\text{civic}} > 0,
\quad
D_{\text{industrial}} \text{ remains high},
\]

but dimensional inequality is reduced:

\[
\Delta D_{\text{industrial-civic}} \downarrow.
\]

The total control dimensionality is distributed rather than centralized.

\subsection{Optimization Pressure}

Let $C(t)$ represent aggregate optimization capability over time.
As automation and computational capacity increase,

\[
\frac{dC}{dt} > 0.
\]

In opaque systems, $D_{\text{civic}}$ remains approximately constant,
while $C(t)$ increases.

Thus, the instability pressure

\[
\Pi(t) =
\left( \frac{C(t)}{D_{\text{civic}}} \right)
\left( \frac{T_{\text{human}}}{\tau_{\text{system}}} \right)
\]

monotonically increases.

In hyperpleonastic systems, $D_{\text{civic}}$ grows as selector-field
artifacts accumulate and RID thresholds are enforced.

Thus, the ratio $C(t)/D_{\text{civic}}$ grows more slowly.

\subsection{Failure Topology}

Opaque economies exhibit centralized failure topology.

Let $k$ represent the number of shared critical dependencies
(e.g., cloud servers, proprietary firmware ecosystems).

Failure probability $P_{\text{collapse}}$ scales superlinearly with $k$:

\[
P_{\text{collapse}} \sim f(k), \quad f''(k) > 0.
\]

Hyperpleonastic economies reduce shared interpretive chokepoints.
Template grammars may remain common, but selector-field recoverability
reduces dependency on synchronized central keys.

Failure probability scales more linearly:

\[
P_{\text{collapse}} \sim g(k), \quad g''(k) \approx 0.
\]

\subsection{Economic Value Migration}

In opaque regimes, economic value concentrates in mechanisms that
restrict interpretive access. Intellectual property secrecy, replacement
part monopolies, firmware licensing, and cloud-mediated diagnostics
become primary sites of rent extraction. Revenue derives less from the
intrinsic performance of artifacts and more from control over the keys
required to interpret, repair, or modify them. Generative grammars are
treated as scarce institutional assets rather than as partially
recoverable structures.

In hyperpleonastic regimes, value migrates toward different domains.
Template innovation, material science advancement, open grammar
maintenance, and tooling precision become the principal competitive
frontiers. Because interpretive authority is distributed across the
material substrate, monopoly power over decoding diminishes. Firms
compete on the quality, efficiency, and durability of their templates
rather than on the opacity of their control systems.

The monopoly of interpretive authority consequently weakens. Economic
advantage shifts from the enclosure of grammar to the refinement of
generative design.

\subsection{The Black-Box Tax as Externality}

Recall the opacity penalty $\Phi$. In opaque economies, artifacts with
high $\Phi$ generate systemic risk externalities. Delayed repair cycles,
supply-chain brittleness, and elevated crisis amplification follow from
the externalization of generative grammars. When interpretive keys are
concentrated and coordination channels are required for restoration,
localized disturbances propagate more readily through the network.

These costs are not borne solely by manufacturers. They are distributed
across the civic substrate, manifesting as prolonged downtime,
substitution scarcity, and collective vulnerability during disruption.
Opacity therefore functions as a negative externality whose burden is
socialized even when the gains from enclosure are privatized.

Hyperpleonastic design internalizes part of this externality by
increasing Recoverable Information Density and distributing interpretive
authority. By embedding partial generative grammar within artifacts
themselves, it reduces dependence on centralized coordination and lowers
the systemic amplification of failure.

\subsection{Stability Curves}

We define civilizational stability margin $\Sigma$ as:

\[
\Sigma = \frac{D_{\text{civic}}}{C}.
\]

In opaque economies:

\[
\frac{d\Sigma}{dt} < 0.
\]

In hyperpleonastic economies:

\[
\frac{d\Sigma}{dt} \approx 0 \quad \text{or} \quad > 0,
\]

depending on enforcement of RID thresholds.

The distinction lies not in technological regression,
but in geometric embedding of legibility.

\section{Slack, Noise, and Agency}

Throughout this essay, tolerance manifolds have been treated as
latent channel capacity and selector fields as mappings from slack
into interpretive structure. We now return to the foundational
assumption underlying this construction: that variation within
functional bounds is not waste, but potential.

\subsection{The Suppression of Slack}

Industrial optimization exhibits a persistent tendency toward
uniformity. Surface finishes are polished to eliminate irregularity,
layer thicknesses are standardized to reduce variance, and deviations
from canonical form are treated as defects. The manufacturing process
is structured to minimize dispersion around target parameters, thereby
reducing uncertainty in quality control and streamlining mass
production.

From a narrow efficiency perspective, slack appears indistinguishable
from noise. Variation that does not directly contribute to immediate
performance is regarded as surplus and therefore as a candidate for
elimination. In this regime, the tolerance manifold is progressively
contracted. Formally, one may represent this contraction as

\[
\mathcal{T} \rightarrow \{ \text{canonical configuration} \},
\]

indicating the collapse of admissible variation into a single optimized
state. Such contraction reduces entropy in production processes,
simplifies logistics, and enhances reproducibility.

However, this same contraction reduces recoverable information. When
variation is eliminated, the channel capacity of the tolerance manifold
declines. Consequently,

\[
I_{\text{recoverable}} \downarrow,
\quad
\text{RID} \downarrow,
\quad
D_{\text{control}} \downarrow.
\]

The artifact may become mechanically refined, yet interpretively
impoverished. Its geometry no longer encodes structured variation
capable of signaling template logic or operational constraints. Smooth
surfaces replace textured grammars; canonical form replaces distributed
memory.

The result is an object that is optimized for production but
attenuated in legibility. It is precise, but silent.

\subsection{Noise as Memory}

Selector-field design reinterprets bounded variation within a tolerance
manifold as a distributed memory substrate. Features that would
otherwise be dismissed as incidental surface irregularities are
structured so as to participate in the encoding of generative grammar.
Pigment mottling, the angular displacement of a ring, the pitch of a
braid, or the visible routing of a heating trace cease to function as
merely aesthetic or manufacturing byproducts. Instead, they become
physically instantiated parameters within an admissible configuration
space.

Within conventional optimization regimes, such variation is treated as
noise to be minimized. Manufacturing seeks canonical uniformity,
reducing deviation to preserve reproducibility and streamline quality
control. Hyperpleonastic design adopts the opposite orientation: it
preserves bounded variation precisely because that variation can encode
template selection, operational limits, or material lineage. The
tolerance manifold becomes a low-amplitude information channel embedded
within functional constraints.

Noise becomes signal when it is structured relative to a decoding
grammar. The decisive transformation is not the magnitude of variation
but its systematic correlation with interpretable states. When
geometric or chromatic variation is coordinated with a template index,
the artifact acquires a distributed memory of its own generative logic.
This memory is not abstract; it is materially instantiated.

Crucially, the signal remains inseparable from the substrate that
carries it. Because encoding occurs within admissible physical
variation, the interpretive structure cannot be revoked without
destroying the artifact itself or collapsing its functional envelope.
The memory of the object is therefore coupled to its persistence.
Legibility endures so long as geometric residue remains.

\subsection{Agency Within the Tolerance Manifold}

Agency requires degrees of freedom.

In a fully compressed system, where variation is eliminated
and interpretive keys are externalized, local agency contracts.

Formally, if
\[
\dim(\mathcal{T}) \rightarrow 0,
\]
then the space of locally selectable interpretations vanishes.

Selector fields preserve a nontrivial $\dim(\mathcal{T})$
and map it into public template grammars.

This mapping enlarges the effective dimensional control surface.

Slack is therefore not inefficiency.
It is the geometric precondition for distributed agency.

\subsection{Entropy and Interpretation}

Entropy increases under degradation.
Materials wear.
Infrastructure fails.

If interpretive grammars are externalized,
legibility may vanish before functionality does.

Selector-field artifacts resist this interpretive entropy.
Their recoverable information density decays slowly:

\[
\frac{d}{dd} \text{RID}(d) \approx 0 \quad \text{for small } d.
\]

The object continues to explain itself
even as it ages.

Robustness becomes the persistence of meaning under stress.

\subsection{Alignment Reconsidered}

Alignment is often framed as a problem of
aligning objectives between humans and machines.

The present analysis proposes a complementary constraint:

\medskip

\noindent
\textbf{Civilizational Alignment Constraint.}
Optimization must not collapse the tolerance manifold
to the point that recoverable information density
falls below the threshold required for distributed control.

\medskip

This is not a rejection of optimization.
It is a geometric limit on its compressive impulse.

\subsection{The Refusal to Become Mute}

An opaque artifact is mute in a precise structural sense. Its
generative grammar exists outside its material substrate, and its
meaning is conditional upon access to external authorization channels.
Interpretation requires permission, and repair requires coordination
beyond the artifact itself. The object functions, but it does not
participate in its own explanation.

A hyperpleonastic artifact, by contrast, preserves a degree of
self-articulation. Its geometry advertises its template class, its
bounded variation encodes selector logic, and its degradation exposes
rather than conceals structural relationships. Damage does not erase
meaning; it transforms the configuration space through which meaning is
recovered.

Incompressibility, in this sense, is not synonymous with gratuitous
complexity. It denotes the preservation of sufficient structured
variation to prevent collapse into a minimal external description.
A materially incompressible system resists reduction to a single
institutional key. It refuses to allow the material world to become
its own simplest description.

\subsection{Conclusion}

Recoverable Information Density, Dimensional Poverty, and the Selector
Field Theorem together describe a unified structural phenomenon. They
articulate how interpretive authority is distributed, how generative
grammars are either embedded or displaced, and how the dimensional
control surface available to agents expands or contracts in response.

The stability of a civilization depends not solely on the intelligence
of its computational systems but on the persistence of legibility within
its material substrate. When artifacts retain structured slack within
their tolerance manifolds, agency remains locally anchored. When that
slack is suppressed, interpretive authority migrates outward, and
fragility increases.

Slack is the geometric reserve in which agency resides. Noise, when
structured relative to a decoding grammar, becomes a substrate of
memory. Selector fields operationalize this principle by embedding
template selection within admissible variation. They constitute a
geometry of trust grounded not in centralized verification but in
material recoverability.

To preserve dimensional control is to preserve the capacity for repair,
reinterpretation, and local autonomy. Civilizational incompressibility
is therefore not an appeal to inefficiency or romantic nostalgia. It is
the maintenance of a stability margin in a regime of accelerating
optimization, ensuring that generative grammars remain partially
recoverable even as production processes converge toward canonical form.


\newpage
\section*{Appendices}

\appendix

\section{Algorithmic Information and Generative Compression}

The main text employed Shannon entropy and mutual information
to define Recoverable Information Density (RID). While Shannon
information captures distributional uncertainty, it does not
directly measure generative compressibility.

We therefore introduce a complementary algorithmic formulation.

\subsection{Kolmogorov Complexity of Generative Grammar}

Let $\mathcal{G}$ denote the generative grammar of an artifact.
The Kolmogorov complexity of $\mathcal{G}$ is defined as

\[
K(\mathcal{G}) = \min_{p} \{ |p| : U(p) = \mathcal{G} \},
\]

where $U$ is a universal Turing machine and $|p|$ is the length
of the shortest program that produces $\mathcal{G}$.

$K(\mathcal{G})$ measures the minimal description length of the
artifact’s generative logic.

An artifact is algorithmically compressible if

\[
K(\mathcal{G}) \ll |\mathcal{G}|,
\]

where $|\mathcal{G}|$ denotes the explicit description length.

\subsection{Externalized Compression}

In opaque systems, the shortest description of the artifact’s
functional logic is not recoverable from the artifact itself.

Formally, if $A$ denotes the physical artifact, then

\[
K(\mathcal{G} \mid A) \approx K(\mathcal{G}),
\]

meaning the artifact provides negligible compression of its
own generative description.

The minimal program reconstructing $\mathcal{G}$ must be obtained
externally.

This condition corresponds to low RID and high opacity penalty $\Phi$.

\subsection{Selector Fields as Partial Programs}

In hyperpleonastic systems, selector fields embed partial programs
within the tolerance manifold.

Let $S(A)$ denote the selector field recoverable from the artifact.
Then the conditional Kolmogorov complexity becomes

\[
K(\mathcal{G} \mid A) = K(\mathcal{G} \mid S(A)) \ll K(\mathcal{G}).
\]

The artifact itself reduces the description length required to
reconstruct its grammar.

Selector fields do not necessarily encode the full program,
but they encode an index into a template library,
thereby shortening the effective program length.

\subsection{Algorithmic Incompressibility and Civic Stability}

We may now reinterpret civilizational incompressibility.

A system is algorithmically centralized if the shortest program
describing its generative grammars resides within a small
institutional subset of agents.

It is algorithmically distributed if the conditional complexity

\[
K(\mathcal{G} \mid A)
\]

is uniformly bounded across artifacts accessible to citizens.

Dimensional Poverty corresponds to a regime in which

\[
K(\mathcal{G} \mid A) \approx K(\mathcal{G})
\]

for most critical artifacts.

Civic stability requires that for essential systems,

\[
K(\mathcal{G} \mid A) \leq K(\mathcal{G}) - \Delta,
\]

for some nontrivial $\Delta > 0$,
ensuring that artifacts materially encode portions of their own
shortest descriptions.

\subsection{Compression, Optimization, and Collapse}

Optimization often reduces surface variance and enforces canonical
forms. While this may reduce manufacturing entropy, it can increase
algorithmic centralization.

If the minimal description of generative grammar is fully externalized,
then systemic failure or institutional collapse renders the grammar
irrecoverable.

Algorithmic incompressibility at the civic level does not imply
inefficiency. It implies distributed possession of partial programs
sufficient for reconstruction.

Selector fields therefore function as embedded description-length
reducers, preserving algorithmic recoverability under bounded
infrastructure failure.

\section{Minimum Description Length and Template Libraries}

The previous appendix introduced Kolmogorov complexity as a
measure of generative compressibility. We now refine the analysis
using the Minimum Description Length (MDL) framework, which
formalizes model selection under finite data.

\subsection{Template Libraries as Model Classes}

Let $\{\mathcal{M}_i\}_{i \in \mathcal{I}}$ denote a finite
template library, where each $\mathcal{M}_i$ represents a
parameterized generative grammar for a class of artifacts.

An artifact $A$ is generated by some $\mathcal{M}_{i^*}$ with
parameters $\theta^*$.

Under MDL, the total description length of $A$ is

\[
L(A) = L(i^*) + L(\theta^* \mid i^*) + L(A \mid i^*, \theta^*),
\]

where:

\begin{itemize}
\item $L(i^*)$ encodes the model index,
\item $L(\theta^* \mid i^*)$ encodes model parameters,
\item $L(A \mid i^*, \theta^*)$ encodes residual variation.
\end{itemize}

In opaque systems, $L(i^*)$ is not recoverable from the artifact.
The model index resides externally.

Thus, reconstruction requires full external model retrieval.

\subsection{Selector Fields as Model Index Encoders}

In hyperpleonastic systems, the selector field $S(A)$ encodes $i^*$.

Formally,

\[
S(A) = i^*.
\]

This reduces effective description length required for reconstruction:

\[
L(A \mid S(A)) = L(\theta^* \mid i^*) + L(A \mid i^*, \theta^*).
\]

The artifact itself supplies the model index.
External search over the model class is eliminated.

Selector fields therefore act as MDL-optimal model selectors
embedded within the tolerance manifold.

\subsection{Tolerance Manifold and Channel Capacity}

Let $\mathcal{T}$ denote the tolerance manifold of dimension $d$.
Suppose variation along each dimension can be discretized into
$b$ admissible bins without violating functionality.

The maximum template index capacity is bounded by

\[
|\mathcal{I}| \leq b^d.
\]

Thus, the channel capacity of the tolerance manifold satisfies

\[
C_{\mathcal{T}} \leq d \log_2 b.
\]

Collapsing $\mathcal{T}$ to a canonical configuration sets $d=0$,
and therefore

\[
C_{\mathcal{T}} = 0.
\]

No template index can be embedded.

Preserving slack preserves channel capacity.

\subsection{Selector-Field Efficiency}

An ideal selector field maximizes mutual information

\[
I(S(A); i^*)
\]

subject to the constraint that perturbations remain within
functional bounds.

Over-encoding that interferes with mechanical integrity
violates the functional envelope.
Under-encoding leaves channel capacity unused.

The design problem is therefore constrained optimization:

\[
\max_{S} \ I(S(A); i^*)
\quad
\text{subject to}
\quad
A \in \mathcal{F}.
\]

This formalizes hyperpleonastic design as
an information-constrained optimization within physical limits.

\subsection{Distributed Model Stability}

If model indices are embedded locally,
then model selection becomes decentralized.

Let $\mathcal{H}$ denote a set of households.
If each artifact encodes its own $i^*$,
then model recovery does not require global search.

The systemic description length of the economy reduces from

\[
L_{\text{centralized}} = \sum_{i \in \mathcal{I}} L(i)
\]

to

\[
L_{\text{distributed}} = \sum_{A \in \mathcal{A}} L(S(A)),
\]

where template indices are redundantly embedded.

Redundancy increases robustness under partial loss.

\subsection{Conclusion of Technical Formalism}

Under the MDL interpretation, selector fields embed
model-selection instructions within admissible geometric variation.

Tolerance manifold dimension determines encoding capacity.
Recoverable Information Density measures practical inference capacity.
Dimensional control surface measures intervention capacity.

Civilizational incompressibility corresponds to the
distributed availability of model indices sufficient
to reconstruct essential generative grammars.

\section{Categorical Formulation of Selector Fields}

The previous appendices treated selector fields as mappings
between sets and model indices. We now refine this structure
categorically.

The purpose is not abstraction for its own sake,
but to clarify how selector fields preserve structure
between admissible geometric variation and template grammars.

\subsection{The Category of Tolerance Configurations}

Let $\mathbf{Tol}$ denote the category whose objects are admissible
configurations $t \in \mathcal{T}$ contained within the functional
envelope $\mathcal{F}$. These objects represent physically realizable
states of the artifact that remain compatible with its operational
constraints. Morphisms in $\mathbf{Tol}$ are physically admissible
deformations $f : t \to t'$ that preserve functionality, mapping one
valid configuration to another without violating the functional
envelope.

Composition in $\mathbf{Tol}$ corresponds to successive admissible
deformations. If $f : t \to t'$ and $g : t' \to t''$ are morphisms,
then $g \circ f : t \to t''$ represents the cumulative transformation
obtained by applying $f$ followed by $g$. Identity morphisms correspond
to null deformations, representing the persistence of a configuration
under no change.

Thus, $\mathbf{Tol}$ formalizes the structured space of permissible
variation, encoding not only the configurations themselves but the
allowable transitions between them.

Thus, $\mathbf{Tol}$ encodes the structure of the tolerance manifold,
including how configurations can transform without violating function.

\subsection{The Category of Templates}

Let $\mathbf{Temp}$ denote the category whose:

\begin{itemize}
\item Objects are template grammars $\mathcal{M}_i$.
\item Morphisms are refinement or specialization maps
      $\phi : \mathcal{M}_i \to \mathcal{M}_j$
      preserving generative compatibility.
\end{itemize}

For example, a base material template may refine into
a batch-specific specialization.

Composition corresponds to successive refinement.

\subsection{Selector Fields as Functors}

A selector field is not merely a function
$S : \mathcal{T} \to \mathcal{I}$.
It is a functor:

\[
S : \mathbf{Tol} \to \mathbf{Temp}.
\]

This requires:

\begin{itemize}
\item For each object $t \in \mathbf{Tol}$,
      $S(t) = \mathcal{M}_i$ for some template.
\item For each morphism $f : t \to t'$,
      there exists a morphism
      $S(f) : S(t) \to S(t')$
      in $\mathbf{Temp}$.
\end{itemize}

Thus, admissible geometric deformation
induces admissible template refinement.

The selector field preserves structure.

\subsection{Functorial Stability}

The functor $S$ must satisfy:

\[
S(\text{id}_t) = \text{id}_{S(t)},
\quad
S(g \circ f) = S(g) \circ S(f).
\]

This ensures interpretive consistency.
Small physical perturbations correspond to
coherent template refinements.

Non-functorial mappings would produce
interpretive discontinuities under minor deformation,
violating graceful degradation.

\subsection{Recoverability as Faithfulness}

Recoverable Information Density corresponds categorically
to partial faithfulness of the functor $S$.

If $S$ is not faithful,
distinct tolerance configurations collapse
to indistinguishable template objects.

If $S$ is faithful on a subcategory
of critical configurations,
then structural distinctions remain observable.

Dimensional Poverty corresponds to a degenerate functor
that factors through a terminal object:

\[
S : \mathbf{Tol} \to \mathbf{1}.
\]

All configurations map to a single opaque template.
Interpretive variation vanishes.

\subsection{Distributed Grammars as Colimits}

Let $\{\mathcal{M}_i\}$ form a diagram in $\mathbf{Temp}$.
The global template grammar may be expressed as a colimit:

\[
\mathcal{M}_{\text{global}} = \varinjlim \mathcal{M}_i.
\]

Selector-field artifacts embed indices to local objects
within this diagram.

Under partial infrastructure failure,
the colimit may not be fully accessible.
However, local objects $\mathcal{M}_i$ remain recoverable
via the functor $S$.

Distributed legibility thus corresponds to
local availability of objects in $\mathbf{Temp}$,
independent of global colimit reconstruction.

\subsection{Alignment as Natural Transformation Constraint}

Suppose optimization processes act as endofunctors:

\[
F : \mathbf{Tol} \to \mathbf{Tol}.
\]

Optimization seeks to contract variation,
often collapsing morphisms in $\mathbf{Tol}$.

Alignment requires that $S$ remain compatible with $F$,
i.e., there exists a natural transformation:

\[
\eta : S \circ F \Rightarrow S.
\]

This expresses that optimization must not
destroy interpretive structure.

If no such natural transformation exists,
optimization collapses selector fidelity,
and RID decreases.

\subsection{Categorical Interpretation of Incompressibility}

Civilizational incompressibility corresponds to
the preservation of nontrivial functorial mappings
from tolerance categories to template categories.

Compression that reduces $\mathbf{Tol}$
to a single object trivializes the functor $S$.

Incompressibility preserves categorical richness,
ensuring that artifacts continue to participate
in meaningful structural mappings.

Selector fields are therefore not decorative encodings.
They are structure-preserving functors
between geometry and grammar.

\section{Sheaf-Theoretic Interpretation: Holographic Steganography Across Scales}

The categorical appendix framed selector fields as functors preserving
structure between tolerance configurations and template grammars. We now
add a sheaf-theoretic layer, motivated not by civic topology but by the
technical requirement of \emph{holographic steganography}: instructions
repeated across scales, including bootstrapping instructions for building
an interpreter device.

The key idea is that interpretive meaning is \emph{locally recoverable}
from partial observations and \emph{globally consistent} when such local
recoveries are glued.

\subsection{Observation Patches and Multiscale Coverings}

Let $X$ denote the physical substrate of an artifact: its surface,
interior microstructure, weave, braid, or printed layers. We treat $X$
as a topological space of observables, not as a civic environment.

An observer does not access $X$ all at once. They access partial regions:
a visible patch on the surface, a cross-section, a corner seam, or a
segment of weave. Let $\{U_\alpha\}_{\alpha \in A}$ be an open cover of
$X$, where each $U_\alpha$ corresponds to an accessible observation
patch.

To model multiscale repetition, we consider a nested family of covers
indexed by scale $\ell$:
\[
\mathcal{U}^{(1)} \prec \mathcal{U}^{(2)} \prec \cdots
\]
where $\mathcal{U}^{(1)}$ is coarse (macroscopic features) and
$\mathcal{U}^{(k)}$ becomes progressively finer (microtextures, fiber
orientation, layer banding).

\subsection{The Sheaf of Instructions}

Define a presheaf $\mathscr{I}$ on $X$ assigning to each open set $U$ a
set (or structured space) of \emph{instruction fragments} observable on
that region:
\[
\mathscr{I}(U) = \{\text{all decodings of instructions recoverable from } U\}.
\]
Restriction maps
\[
\rho_{UV} : \mathscr{I}(U) \to \mathscr{I}(V), \qquad V \subseteq U,
\]
formalize that a decoding available from a larger patch restricts to a
decoding on a smaller patch.

We impose the sheaf condition: if instruction fragments agree on
overlaps, they glue to a unique global instruction.

\medskip
\noindent
\textbf{Sheaf Condition.}
Given an open cover $\{U_\alpha\}$ of $U$, a family
$s_\alpha \in \mathscr{I}(U_\alpha)$ that satisfies
\[
\rho_{U_\alpha, U_\alpha \cap U_\beta}(s_\alpha)
=
\rho_{U_\beta, U_\alpha \cap U_\beta}(s_\beta)
\quad \text{for all } \alpha,\beta
\]
glues to a unique $s \in \mathscr{I}(U)$ such that
$\rho_{U,U_\alpha}(s)=s_\alpha$ for all $\alpha$.

\medskip

This expresses holographic recoverability: local readings can be glued
to consistent global interpretation, and inconsistencies indicate either
damage or adversarial tampering.

\subsection{Redundancy as Multiscale Sections}

Holographic steganography requires that meaningful instruction sections
exist at multiple scales.

We therefore posit that for many $U$ in coarse covers
$\mathcal{U}^{(1)}$, $\mathscr{I}(U)$ already contains nontrivial
sections corresponding to high-level classification and safe operation,
while finer covers provide more specific sections: parameters, batch
lineage, repair pathways.

Formally, we model a filtration of subsheaves
\[
\mathscr{I}_{\text{coarse}} \subseteq
\mathscr{I}_{\text{mid}} \subseteq
\mathscr{I}_{\text{fine}} \subseteq \mathscr{I},
\]
where inclusion corresponds to increasing interpretive resolution.

This matches the layered interpretability requirement in the main text,
but now expressed as a sheaf-theoretic refinement of sections.

\subsection{The Sheaf of Interpreters}

To avoid dependence on a privileged decoding apparatus, the artifact
must include instructions not only for the target system but also for
constructing an interpreter device.

We formalize this by introducing a second sheaf $\mathscr{D}$ assigning
to each open set $U$ the space of \emph{interpreter descriptions}
recoverable from $U$:
\[
\mathscr{D}(U) = \{\text{device schematics / decoding protocols inferable from } U\}.
\]

The essential bootstrapping requirement is that there exist nontrivial
sections of $\mathscr{D}$ on coarse patches:
\[
\exists\, U \in \mathcal{U}^{(1)} \text{ such that } \mathscr{D}(U) \neq \varnothing.
\]
In plain terms: macroscopic features must suffice to begin constructing a
decoder, even if that decoder is crude.

Finer patches then refine the decoder: calibration, error correction,
measurement sensitivity, and additional template families.

Thus, interpretation is not a single sheaf but a coupled pair:
\[
(\mathscr{I}, \mathscr{D}).
\]

\subsection{Bootstrapping as Mutual Gluing}

The instruction sheaf $\mathscr{I}$ and interpreter sheaf $\mathscr{D}$
are mutually constraining.

A decoder section $d \in \mathscr{D}(U)$ induces an evaluation map on
instruction fragments:
\[
\operatorname{eval}_d : \mathscr{I}(U) \to \text{DecodedMeaning}(U).
\]

Conversely, instruction fragments may contain calibration constraints
that restrict admissible decoders. This suggests a compatibility
relation, which can be formalized as a sheaf morphism:
\[
\mathscr{D} \to \underline{\operatorname{Hom}}(\mathscr{I}, \mathscr{M}),
\]
where $\mathscr{M}$ is a sheaf of meanings (e.g., template indices,
parameter sets, repair procedures) and
$\underline{\operatorname{Hom}}$ denotes an internal hom object when
$\mathscr{I}$ and $\mathscr{M}$ take values in a suitable category.

In practical terms: the artifact supplies both a text and a key,
but distributed as local sections that glue.

\subsection{Damage, Ambiguity, and Cohomological Obstructions}

When an artifact is damaged, local readings may fail to agree on
overlaps. In the sheaf formalism, such failures appear as obstructions
to gluing.

Let $s_\alpha \in \mathscr{I}(U_\alpha)$ be local decodings on a cover.
If agreement fails on overlaps, then the family $\{s_\alpha\}$ defines a
nontrivial \v{C}ech cocycle measuring inconsistency.

This suggests an operational interpretation: repair is the process of
modifying the substrate until the cocycle becomes trivial and the
instructions glue.

Thus, the sheaf formalism provides an intrinsic error-detection and
tamper-detection criterion: consistency on overlaps is the certificate of
integrity.

\subsection{Holographic Steganography Principle}

We may state the core design requirement as follows.

\medskip
\noindent
\textbf{Holographic Steganography Principle.}
An artifact implements hyperpleonastic legibility if there exist sheaves
$(\mathscr{I},\mathscr{D})$ on its physical substrate $X$ such that:
(i) nontrivial instruction sections and nontrivial interpreter sections
exist at coarse observational scale; (ii) finer observational scales
provide refinements that glue consistently; and (iii) damage manifests as
explicit gluing obstructions detectable from local overlap checks.

\medskip

This principle captures the intended ``instructions repeated across
scales'' requirement, including self-hosting instructions for constructing
an interpreter device, while remaining independent of any particular
material technology.

\newpage
\begin{thebibliography}{99}

\bibitem{shannon1948}
C. E. Shannon,
``A Mathematical Theory of Communication,''
\emph{Bell System Technical Journal},
vol. 27, pp. 379--423, 623--656, 1948.

\bibitem{ashby1956}
W. R. Ashby,
\emph{An Introduction to Cybernetics},
Chapman \& Hall, 1956.

\bibitem{wiener1948}
N. Wiener,
\emph{Cybernetics: Or Control and Communication in the Animal and the Machine},
MIT Press, 1948.

\bibitem{nyquist1932}
H. Nyquist,
``Regeneration Theory,''
\emph{Bell System Technical Journal},
vol. 11, pp. 126--147, 1932.

\bibitem{yudkowsky2008}
E. Yudkowsky,
``Artificial Intelligence as a Positive and Negative Factor in Global Risk,''
in \emph{Global Catastrophic Risks},
Bostrom and Cirkovic (eds.),
Oxford University Press, 2008.

\bibitem{yampolskiy2015}
R. V. Yampolskiy,
\emph{Artificial Superintelligence: A Futuristic Approach},
CRC Press, 2015.

\bibitem{bostrom2014}
N. Bostrom,
\emph{Superintelligence: Paths, Dangers, Strategies},
Oxford University Press, 2014.

\bibitem{doctorow2023}
C. Doctorow,
\emph{The Internet Con: How to Seize the Means of Computation},
Verso, 2023.

\bibitem{perens1999}
B. Perens,
``The Open Source Definition,''
Open Source Initiative, 1999.

\bibitem{gallagher2017}
S. Gallagher,
``The Right to Repair Movement,''
\emph{Ars Technica},
2017.

\bibitem{jacobson1995}
T. Jacobson,
``Thermodynamics of Spacetime: The Einstein Equation of State,''
\emph{Physical Review Letters},
vol. 75, no. 7, 1995.

\bibitem{latour1992}
B. Latour,
``Where Are the Missing Masses? The Sociology of a Few Mundane Artifacts,''
in \emph{Shaping Technology / Building Society},
MIT Press, 1992.

\bibitem{holling1973}
C. S. Holling,
``Resilience and Stability of Ecological Systems,''
\emph{Annual Review of Ecology and Systematics},
vol. 4, 1973.

\bibitem{ostrom1990}
E. Ostrom,
\emph{Governing the Commons},
Cambridge University Press, 1990.

\bibitem{simon1962}
H. A. Simon,
``The Architecture of Complexity,''
\emph{Proceedings of the American Philosophical Society},
vol. 106, no. 6, 1962.

\end{thebibliography}

\end{document}
